\documentclass[10pt,a4paper]{amsart}
\usepackage[margin=19mm]{geometry}
\usepackage{amsmath,amssymb,mathtools}
\usepackage[scr=rsfs,cal=euler]{mathalfa}
\usepackage{xfrac}
\usepackage[unicode,hidelinks,bookmarks=false]{hyperref}
\newcommand{\ie}{i.e.\ }
\newcommand{\cf}{cf.\ }
\newcommand{\resp}{respectively\ }
\newcommand{\Subsection}{\subsection}
\providecommand{\nobreakdash}{\nobreak\hbox{-}}
\setlength{\emergencystretch}{2em}
\multlinegap0pt
\usepackage{bm}
\usepackage{bbm}
\usepackage{dsfont}
\usepackage{xspace}
\usepackage{cancel}
\usepackage{mathtools}
\usepackage{cleveref}
\usepackage[shortlabels]{enumitem}
\usepackage{amsthm}
\makeatletter
\@ifundefined{theorem}{\newtheorem{theorem}{Theorem}}{}
\@ifundefined{lemma}{\newtheorem{lemma}[theorem]{Lemma}}{}
\makeatother
\newcommand{\CNOT}{\mathrm{CNOT}}
\newcommand{\SU}{\mathrm{SU}}
\newcommand{\bbC}{\mathbb{C}}
\newcommand{\bbR}{\mathbb{R}}
\newcommand{\cM}{\mathcal{M}}
\title[Identity-Paired Progressive Depth Training: When Trainabilit]{Identity-Paired Progressive Depth Training: When Trainability Persists Beyond Expressibility}
\author{Athanasios Hadjidimoulas, Tirthak Patel,, Anastasios Kyrillidis}
\date{Source: arXiv:2607.16800v2 (CC BY 4.0)}
\begin{document}
\maketitle

\noindent\emph{Source and licence.} Identity-Paired Progressive Depth Training: When Trainability Persists Beyond Expressibility. Version arXiv:2607.16800v2 (CC BY 4.0), distributed under the Creative Commons Attribution 4.0 International licence, as recorded in the arXiv licence metadata for this exact version and shown on the abstract page https://arxiv.org/abs/2607.16800v2. Full rights notice: licenses/arxiv-2607.16800-CC-BY-4.0.txt.\par\smallskip

\noindent\textit{Editorial scope.} Counted unit: the Theorem whose source label is \texttt{thm:saturation} in the article (source0.tex lines 1162--1171, source label \texttt{thm:saturation}), delivered together with the article's own complete original proof of it (source0.tex lines 1172--1372, one \texttt{proof} environment of 201 lines). No mathematics is written, repaired or re-derived: statement and proof are the article's own text line for line. Required context retained uncounted: eq:pauli\_rot (lines 327--332); eq:inv\_rot (lines 346--349); eq:cnot\_ring (lines 458--461); eq:UD (lines 758--763); eq:Wk (lines 765--771); lem:factor (lines 1096--1108); eq:state\_form (lines 1099--1104); lem:su2\_cover (lines 2654--2657). Every definition, display and label of those lines is retained unchanged. Omitted from this package and not redistributed: 1--326, 333--345, 350--457, 462--757, 764--764, 772--1095, 1105--1161, 1373--2653, 2658--4020. Nothing inside a retained excerpt is omitted or altered: every retained body is delivered exactly as the article's lines are written, so no omission receipt of any kind accompanies this package. Citations: the retained excerpts carry no citation command at all, so no bibliography entry of the article is transcribed and no reference is redistributed. Reference adaptation: none was needed and none was made; every reference inside the delivered unit resolves inside it, so no unresolved reference or citation is produced. Printed slips kept literal: no printed word, symbol, hypothesis or number of the counted statement or of its proof was altered. Layout and class: the article's own class is \texttt{article} with packages including geometry, amsmath, amssymb, amsthm, algorithm, algpseudocode, booktabs, hyperref, cleveref, graphicx, xcolor, enumitem, bm, mathtools; the wrapper is the lane's pinned \texttt{amsart} editorial wrapper at 10pt on A4, which restates the theorem-like environments the retained text needs and adds the article's own macro definitions that the retained text needs (\texttt{\textbackslash CNOT}, \texttt{\textbackslash SU}, \texttt{\textbackslash bbC}, \texttt{\textbackslash bbR}, \texttt{\textbackslash cM}), with extra wrapper packages (bm, bbm, dsfont, xspace, cancel, mathtools, cleveref, enumitem). The delivered numbering is the unit's own. Assets: no retained span contains a figure environment, a graphics-inclusion command, a PDF-page inclusion, a table or a TikZ picture; no image file is copied into this package, the package declares no asset dependency and no image file is redistributed with it. Licence: version arXiv:2607.16800v2 of this article is distributed under the Creative Commons Attribution 4.0 International licence, as recorded in the arXiv licence metadata for this exact version and shown on the abstract page https://arxiv.org/abs/2607.16800v2. A full rights notice is delivered at licenses/arxiv-2607.16800-CC-BY-4.0.txt. Characters: every retained excerpt is pure ASCII, so the delivered engine, XeLaTeX with its default Latin Modern text font, reports no missing character.
\begin{equation}\label{eq:pauli_rot}
  \mathbf{R}_{\sigma}(\theta)
  = \exp \bigl(-i\tfrac{\theta}{2}\sigma\bigr)
  = \cos \bigl(\tfrac{\theta}{2}\bigr) \mathbf{I}_2
    - i\sin \bigl(\tfrac{\theta}{2}\bigr) \sigma.
\end{equation}

\begin{equation}\label{eq:inv_rot}
  \mathbf{R}^{\mathrm{inv}}(\alpha,\beta,\gamma)
  = \mathbf{R}_X(-\gamma) \mathbf{R}_Z(-\beta) \mathbf{R}_Y(-\alpha).
\end{equation}

\begin{equation}\label{eq:cnot_ring}
  \mathbf{U}_{\mathrm{ent}}
  =\prod_{q=1}^{n}\CNOT_{q\to(q\bmod n)+1},
\end{equation}

\begin{equation}\label{eq:UD}
  \mathbf{U}_D^{\mathrm{SE}}(\bm\theta)
  = \left(\prod_{k=D-1}^{1}\mathbf{W}_k(\bm\theta^{(k)})\right)
    \cdot \mathbf{U}_{\mathrm{ent}}
    \cdot\mathcal{R}^{\mathrm{fwd}}(\bm\theta^{(0)}),
\end{equation}

\begin{equation}\label{eq:Wk}
  \mathbf{W}_k(\bm\theta^{(k)})=
  \begin{cases}
    \mathcal{R}^{\mathrm{inv}}(\bm\theta^{(k)}), & k\text{ odd},\\
    \mathcal{R}^{\mathrm{fwd}}(\bm\theta^{(k)}), & k\ge 2\text{ even}.
  \end{cases}
\end{equation}

\begin{lemma}[State factorization]\label{lem:factor}
  For any depth $D\ge 1$, every state
  $|\boldsymbol{\psi}\rangle\in\cM_D$ can be written as
  \begin{equation}\label{eq:state_form}
    |\boldsymbol{\psi}\rangle
    = \Bigl(\bigotimes_{q=1}^n \mathbf{V}_q\Bigr) 
      \mathbf{U}_{\mathrm{ent}} 
      \Bigl(\bigotimes_{q=1}^n r_q\Bigr)|0^n\rangle,
  \end{equation}
  where $r_q\in\SU(2)$ is the pre-entangler rotation on qubit~$q$, and
  $\mathbf{V}_q = \mathbf{W}_{D{-}1,q}\cdots \mathbf{W}_{1,q}\in\SU(2)$ is the composition of all
  post-entangler rotation layers on qubit~$q$.
\end{lemma}

\begin{lemma}\label{lem:su2_cover}
  The map $\bbR^3\ni(\alpha,\beta,\gamma)\mapsto
  \mathbf{R}_X(-\gamma) \mathbf{R}_Z(-\beta) \mathbf{R}_Y(-\alpha)$ is surjective onto $\SU(2)$.
\end{lemma}

\begin{theorem}[Reachable Set Saturation]\label{thm:saturation}
  For the single-entangler circuit family~\eqref{eq:UD}:
  \begin{enumerate}[label=(\alph*)]
    \item $\cM_1\subseteq\cM_2$.
    \item For all $D\ge 2$, $\cM_D=\cM_2$.
    \item If $\mathbf{U}_{\mathrm{ent}}$ does not normalize the local-unitary
      group $\mathcal{L}=\{\bigotimes_q\mathbf{V}_q:\mathbf{V}_q\in\SU(2)\}$
      (as the CNOT ring does not), then $\cM_1\subsetneq\cM_2$.
  \end{enumerate}
\end{theorem}

\begin{proof}
Throughout, $\cM_D=\{\mathbf{U}_D^{\mathrm{SE}}(\bm\theta)|0^n\rangle:
\bm\theta\in\bbR^{3nD}\}$ is the set of states reachable by the
depth-$D$ SE circuit over all parameter choices.

\textit{Part~(a): $\cM_1\subseteq\cM_2$.}
Let $|\boldsymbol{\psi}\rangle\in\cM_1$, so
$|\boldsymbol{\psi}\rangle=\mathbf{U}_1^{\mathrm{SE}}(\bm\theta^{(0)})|0^n\rangle$ for
some $\bm\theta^{(0)}\in\bbR^{3n}$. By~\eqref{eq:UD} the depth-$1$
circuit has an empty post-entangler product, so
$\mathbf{U}_1^{\mathrm{SE}}(\bm\theta^{(0)})
=\mathbf{U}_{\mathrm{ent}} \mathcal{R}^{\mathrm{fwd}}(\bm\theta^{(0)})$.
Consider the depth-$2$ circuit with parameters
$(\bm\theta^{(0)},\bm 0)$. Its single post-entangler layer is
$\mathbf{W}_1(\bm 0)=\mathcal{R}^{\mathrm{inv}}(\bm 0)$ by~\eqref{eq:Wk}
($k=1$ is odd). Each single-qubit factor of
$\mathcal{R}^{\mathrm{inv}}(\bm 0)$ is
$\mathbf{R}^{\mathrm{inv}}(0,0,0)=\mathbf{R}_X(0)\mathbf{R}_Z(0)\mathbf{R}_Y(0)=\mathbf{I}_2$
by~\eqref{eq:inv_rot} and $\mathbf{R}_\sigma(0)=\mathbf{I}_2$, so
$\mathcal{R}^{\mathrm{inv}}(\bm 0)=\bigotimes_{q=1}^n\mathbf{I}_2=\mathbf{I}_{2^n}$.
Hence
\begin{equation}\label{eq:sat_a}
  \mathbf{U}_2^{\mathrm{SE}}(\bm\theta^{(0)},\bm 0)
  = \mathbf{I}_{2^n} \mathbf{U}_{\mathrm{ent}} 
    \mathcal{R}^{\mathrm{fwd}}(\bm\theta^{(0)})
  = \mathbf{U}_1^{\mathrm{SE}}(\bm\theta^{(0)}),
\end{equation}
so $|\boldsymbol{\psi}\rangle=\mathbf{U}_2^{\mathrm{SE}}(\bm\theta^{(0)},\bm 0)|0^n\rangle
\in\cM_2$. As $|\boldsymbol{\psi}\rangle$ was arbitrary, $\cM_1\subseteq\cM_2$.

\textit{Part~(b): $\cM_D=\cM_2$ for all $D\ge 2$.}
We prove the two inclusions separately.

\emph{($\cM_2\subseteq\cM_D$).} Fix $D\ge 2$ and let
$|\boldsymbol{\psi}\rangle\in\cM_2$ be produced by parameters
$(\bm\theta^{(0)},\bm\theta^{(1)})$. Define depth-$D$ parameters by
keeping $\bm\theta^{(0)},\bm\theta^{(1)}$ and setting
$\bm\theta^{(k)}=\bm 0$ for all $2\le k\le D-1$. For each such $k$,
$\mathbf{W}_k(\bm 0)$ equals $\mathcal{R}^{\mathrm{inv}}(\bm 0)$ or
$\mathcal{R}^{\mathrm{fwd}}(\bm 0)$ by~\eqref{eq:Wk}; either way every
single-qubit factor is $\mathbf{R}^{\mathrm{fwd}}(0,0,0)=\mathbf{I}_2$ or
$\mathbf{R}^{\mathrm{inv}}(0,0,0)=\mathbf{I}_2$, so $\mathbf{W}_k(\bm 0)=\mathbf{I}_{2^n}$. The
post-entangler product therefore collapses to
$\prod_{k=D-1}^{1}\mathbf{W}_k=\mathbf{W}_1(\bm\theta^{(1)})$, giving
$\mathbf{U}_D^{\mathrm{SE}}(\bm\theta^{(0)},\bm\theta^{(1)},\bm 0,\dots,\bm 0)
=\mathbf{U}_2^{\mathrm{SE}}(\bm\theta^{(0)},\bm\theta^{(1)})$ and hence
$|\boldsymbol{\psi}\rangle\in\cM_D$.

\emph{($\cM_D\subseteq\cM_2$).} Let $|\boldsymbol{\psi}\rangle\in\cM_D$ with
$D\ge 2$. By \Cref{lem:factor} it has the form~\eqref{eq:state_form},
\begin{equation}\label{eq:sat_b_form}
  |\boldsymbol{\psi}\rangle
  = \Bigl(\bigotimes_{q=1}^n \mathbf{V}_q\Bigr) \mathbf{U}_{\mathrm{ent}} 
    \Bigl(\bigotimes_{q=1}^n r_q\Bigr)|0^n\rangle,
  \qquad r_q,\mathbf{V}_q\in\SU(2),
\end{equation}
where $\mathbf{V}_q=\mathbf{W}_{D-1,q}\cdots\mathbf{W}_{1,q}$ is the composition of all
post-entangler single-qubit rotations on qubit~$q$ and lies in $\SU(2)$
by closure of $\SU(2)$ under multiplication. At depth $2$ the
post-entangler factor on qubit~$q$ is the single $\SU(2)$ element
$\mathbf{W}_{1,q}=\mathbf{R}^{\mathrm{inv}}(\alpha_q^{(1)},\beta_q^{(1)},\gamma_q^{(1)})$,
and by \Cref{lem:su2_cover} the map
$(\alpha,\beta,\gamma)\mapsto\mathbf{R}^{\mathrm{inv}}(\alpha,\beta,\gamma)$ is
surjective onto $\SU(2)$. Hence for each $q$ there exist angles
$(\alpha_q^{(1)},\beta_q^{(1)},\gamma_q^{(1)})$ with
$\mathbf{R}^{\mathrm{inv}}(\alpha_q^{(1)},\beta_q^{(1)},\gamma_q^{(1)})=\mathbf{V}_q$;
likewise, surjectivity of the forward map---the forward Euler
sequence of \Cref{lem:su2_cover}, Part~(a)---supplies
$\bm\theta^{(0)}$ with
$\mathcal{R}^{\mathrm{fwd}}(\bm\theta^{(0)})=\bigotimes_q r_q$. With this
choice $\mathbf{U}_2^{\mathrm{SE}}(\bm\theta^{(0)},\bm\theta^{(1)})|0^n\rangle$
equals the right-hand side of~\eqref{eq:sat_b_form}, so
$|\boldsymbol{\psi}\rangle\in\cM_2$. The crux is that one post-entangler layer
already attains \emph{every} per-qubit factor $\mathbf{V}_q\in\SU(2)$; the extra
layers present when $D>2$ compose additional $\SU(2)$ elements but,
by closure, cannot leave $\SU(2)$, so they enlarge neither the per-qubit
factor nor $\cM_D$. Combining both inclusions, $\cM_D=\cM_2$.

\textit{Part~(c): $\cM_1\subsetneq\cM_2$ when $\mathbf{U}_{\mathrm{ent}}$ does
not normalize $\mathcal{L}$.}
By Part~(a), $\cM_1\subseteq\cM_2$, so it suffices to exhibit one state
in $\cM_2\setminus\cM_1$. We first give an explicit witness for $n=2$
with $\mathbf{U}_{\mathrm{ent}}=\CNOT_{1\to2}$, the entangling gate from which
the CNOT ring~\eqref{eq:cnot_ring} is built, and then give the general
non-normalization argument. Write
$c_8=\cos(\pi/8)$, $s_8=\sin(\pi/8)$, and recall from~\eqref{eq:pauli_rot}
that $\mathbf{R}_Y(\theta)|0\rangle=\cos\tfrac\theta2|0\rangle
+\sin\tfrac\theta2|1\rangle$ and
$\mathbf{R}_Y(\theta)|1\rangle=-\sin\tfrac\theta2|0\rangle
+\cos\tfrac\theta2|1\rangle$.

\emph{The target depth-$2$ state.} Take the depth-$2$ parameters
realizing $r_1=\mathbf{R}_Y(\pi/2)$, $r_2=\mathbf{I}_2$, $\mathbf{V}_1=\mathbf{I}_2$,
$\mathbf{V}_2=\mathbf{R}_Y(\pi/4)$ in~\eqref{eq:sat_b_form} (each is in $\SU(2)$, hence
attainable by Part~(b)). Applying the gates right-to-left:
\begin{align}
  (\mathbf{R}_Y(\tfrac\pi2)\otimes\mathbf{I}_2)|00\rangle
  &= \bigl(\tfrac{1}{\sqrt2}|0\rangle+\tfrac{1}{\sqrt2}|1\rangle\bigr)
     \otimes|0\rangle
   = \tfrac{1}{\sqrt2}\bigl(|00\rangle+|10\rangle\bigr),
     \label{eq:satc_pre}\\
  \CNOT_{1\to2} \tfrac{1}{\sqrt2}\bigl(|00\rangle+|10\rangle\bigr)
  &= \tfrac{1}{\sqrt2}\bigl(|00\rangle+|11\rangle\bigr),
     \label{eq:satc_ent}\\
  (\mathbf{I}_2\otimes\mathbf{R}_Y(\tfrac\pi4)) 
     \tfrac{1}{\sqrt2}\bigl(|00\rangle+|11\rangle\bigr)
  &= \tfrac{1}{\sqrt2}\bigl(c_8|00\rangle+s_8|01\rangle
     -s_8|10\rangle+c_8|11\rangle\bigr)
   =: |\boldsymbol{\psi}^\star\rangle,
     \label{eq:satc_target}
\end{align}
where~\eqref{eq:satc_ent} uses
$\CNOT_{1\to2}|a\rangle|b\rangle=|a\rangle|a\oplus b\rangle$, and the
second qubit in~\eqref{eq:satc_target} was rotated by
$\mathbf{R}_Y(\pi/4)$ (half-angle $\pi/8$). Thus
$|\boldsymbol{\psi}^\star\rangle\in\cM_2$.

\emph{General form of a $\cM_1$ state.} By~\eqref{eq:UD} with $D=1$,
every $|\boldsymbol{\phi}\rangle\in\cM_1$ is
$\CNOT_{1\to2} \mathcal{R}^{\mathrm{fwd}}(\bm\theta^{(0)})|00\rangle$,
i.e.\ a CNOT applied to a product state. Writing the product state as
$(a|0\rangle+b|1\rangle)\otimes(c|0\rangle+d|1\rangle)$ with
$a,b,c,d\in\bbC$ (the most general single-qubit states; any global
phase is absorbed), expand and apply $\CNOT_{1\to2}$, which fixes
$|00\rangle,|01\rangle$ and swaps the control-$1$ pair
$|10\rangle\leftrightarrow|11\rangle$:
\begin{align}
  (a|0\rangle{+}b|1\rangle)\otimes(c|0\rangle{+}d|1\rangle)
  &= ac |00\rangle+ad |01\rangle+bc |10\rangle+bd |11\rangle,
     \label{eq:satc_prod}\\
  \CNOT_{1\to2}\bigl[ \cdot \bigr]
  &= ac |00\rangle+ad |01\rangle+bd |10\rangle+bc |11\rangle
   =: |\boldsymbol{\phi}\rangle.
     \label{eq:satc_phi}
\end{align}

\emph{The contradiction.} Suppose, for contradiction, that
$|\boldsymbol{\psi}^\star\rangle\in\cM_1$, i.e.\
$|\boldsymbol{\phi}\rangle=|\boldsymbol{\psi}^\star\rangle$ for some $a,b,c,d$. Matching
the four computational-basis amplitudes
of~\eqref{eq:satc_phi} and~\eqref{eq:satc_target} gives
\begin{equation}\label{eq:satc_amps}
  ac=\tfrac{c_8}{\sqrt2},\qquad
  ad=\tfrac{s_8}{\sqrt2},\qquad
  bd=-\tfrac{s_8}{\sqrt2},\qquad
  bc=\tfrac{c_8}{\sqrt2}.
\end{equation}
Since $c_8=\cos(\pi/8)>0$ and $s_8=\sin(\pi/8)>0$, the right-hand sides
$\tfrac{c_8}{\sqrt2}$ and $\pm\tfrac{s_8}{\sqrt2}$ are all nonzero;
hence none of $a,b,c,d$ can vanish (a zero factor would force two of
the products in~\eqref{eq:satc_amps} to be $0$). With $a,b,c,d\neq0$ we
may divide. The first two equations give
\begin{equation}\label{eq:satc_ratio1}
  \frac{ad}{ac}=\frac{d}{c}
  =\frac{s_8/\sqrt2}{c_8/\sqrt2}
  =\tan \Bigl(\frac\pi8\Bigr)>0,
\end{equation}
while the last two give
\begin{equation}\label{eq:satc_ratio2}
  \frac{bd}{bc}=\frac{d}{c}
  =\frac{-s_8/\sqrt2}{c_8/\sqrt2}
  =-\tan \Bigl(\frac\pi8\Bigr)<0.
\end{equation}
Equations~\eqref{eq:satc_ratio1} and~\eqref{eq:satc_ratio2} assign the
single ratio $d/c$ two values of opposite sign,
$+\tan(\pi/8)$ and $-\tan(\pi/8)$, which are distinct because
$\tan(\pi/8)\neq0$. This contradiction shows
$|\boldsymbol{\psi}^\star\rangle\notin\cM_1$, hence
$|\boldsymbol{\psi}^\star\rangle\in\cM_2\setminus\cM_1$ and
$\cM_1\subsetneq\cM_2$.

\emph{General argument (any non-normalizing entangler).}
By the depth-$1$ and depth-$2$ forms,
$\cM_1=\{\mathbf{U}_{\mathrm{ent}}\,|\pi\rangle:|\pi\rangle\ \text{a product state}\}$
and, using the arbitrary post-entangler product layer of Part~(b),
$\cM_2=\{\mathbf{L}\,m:\mathbf{L}\in\mathcal{L},\ m\in\cM_1\}=\mathcal{L}\,\cM_1$,
where $\mathcal{L}$ is the local-unitary group.  Hence $\cM_1=\cM_2$ iff
$\mathcal{L}\,\cM_1=\cM_1$, i.e.\ iff for every $\mathbf{L}\in\mathcal{L}$ and
product $|\pi\rangle$ the state
$\mathbf{U}_{\mathrm{ent}}^{\dagger}\mathbf{L}\,\mathbf{U}_{\mathrm{ent}}|\pi\rangle$ is again a
product state.  The unitaries preserving every product state are exactly
$\mathcal{L}\rtimes S_n$ (local unitaries composed with qubit permutations);
since $\mathcal{L}=\SU(2)^{\otimes n}$ is connected and the conjugation
$\mathbf{L}\mapsto\mathbf{U}_{\mathrm{ent}}^{\dagger}\mathbf{L}\,\mathbf{U}_{\mathrm{ent}}$ fixes
$\mathbf{I}$, continuity confines its image to the identity component
$\mathcal{L}$, so this holds precisely when $\mathbf{U}_{\mathrm{ent}}$
\emph{normalizes} $\mathcal{L}$
($\mathbf{U}_{\mathrm{ent}}^{\dagger}\mathcal{L}\,\mathbf{U}_{\mathrm{ent}}\subseteq\mathcal{L}$).
Contrapositively, if $\mathbf{U}_{\mathrm{ent}}$ does not normalize $\mathcal{L}$
then $\cM_1\subsetneq\cM_2$.  The CNOT ring does not normalize
$\mathcal{L}$: conjugating the local rotation
$\mathbf{R}_X(\alpha)\otimes\mathbf{I}\in\mathcal{L}$ by the entangler gives
$\CNOT_{1\to2}^{\dagger}(\mathbf{R}_X(\alpha)\otimes\mathbf{I})\CNOT_{1\to2}
=e^{-i\alpha\mathbf{X}_1\mathbf{X}_2/2}$, an entangling $XX$-rotation that for
generic $\alpha$ lies outside $\mathcal{L}$, and applying it to a generic
product state produces a member of $\cM_2\setminus\cM_1$.  (By contrast a SWAP
gate, though not a product unitary, \emph{does} normalize $\mathcal{L}$
via $\mathrm{SWAP}\,(\mathbf{u}\otimes\mathbf{v})\,\mathrm{SWAP}=\mathbf{v}\otimes\mathbf{u}$, and
indeed gives $\cM_1=\cM_2$; this is why the correct hypothesis is
non-normalization, not mere non-productness.)
\end{proof}

\end{document}
